\section{Discussion and limitations}
\label{sec:discussion}

\paragraph{What the results say about rule choice.}
First place has two prices, and which rule makes it dearer depends on which binds. If total effort is the constraint, the mean family is the dearest rule on every board (\Cref{sec:pricemap}), because its depth is $m$ times the gap. If effort per task is the constraint, the hardest rule to buy is the minimum on HELM (median minimal feasible budget 11 points, against 6.5 for the mean and 2.7 for plurality) and plurality on the Open LLM Leaderboard (9 and 12 points), and the order moves with the budget (\Cref{sec:breadth}). To make first place costly in total points, average; to demand capability on every task, take the minimum.

\paragraph{Scope.}
The conclusions hold for three boards, 30 top targets, the boost model of \Cref{sec:setting} and budgets seen for fine-tunes of models of 9B and above. Comparisons by total boost are robust to the incumbent on HELM and to the unit when uncapped; those by breadth are not.
\emph{The developer} is a single, static actor: rivals do not respond, so the prices are a lower envelope of the cost on a contested board. \emph{The boost} is an input whose scale is checked, not produced, by training; boosts on different tasks are independent, and a spillover-aware variant lowers the cheapest total boost to 92--95\% of the independent one. \emph{Boards}: HELM is one benchmark family and the Open LLM Leaderboard's top a crowd of derivatives, so each comparison has 11--13 independent clusters.

\paragraph{Impact.}
The prices show how cheaply a board can be gamed; we train nothing and release only an analysis of public scores.
